\uuid{yENV}
\exo7id{938}
\titre{exo7 938}
\auteur{cousquer}
\organisation{exo7}
\datecreate{2003-10-01}
\isIndication{false}
\isCorrection{false}
\chapitre{Application linéaire}
\sousChapitre{Image et noyau, théorème du rang}
\cpge{mpsi-mp2i.applications-lineaires.d, pcsi.applications-lineaires.d, ptsi.applications-lineaires.d, mpsi-mp2i.applications-lineaires.a, pcsi.applications-lineaires.a, ptsi.applications-lineaires.a, bcpst2.applications-lineaires-matrices.b}
\module{Algèbre}
\niveau{L1}
\difficulte{}

\contenu{
\question{
Soit $E$, $F$, $G$ trois espaces vectoriels de dimensions respectives $n$,
$p$, $q$, $f$ et $g$ deux applications linéaires
$E\overset{f}{\rightarrow} F\overset{g}{\rightarrow} G$ telles que
$g\circ f=0$. Quelle relation existe-t-il entre le rang de $f$ et celui de $g$~?
}
}
